% ==========================================================
% titles/title-12.tex -- Title XII: Transitional Provisions
% ==========================================================

\ConstTitle{Transitional Provisions}{title:transitional}

\Article{Founding Mechanisms}{art:founding-mechanisms}

\ArtSection{Qualification Commons Bootstrap}{art:founding-mechanisms:s-qualification-commons-bootstrap}
At founding, before a pool of former \gls{gls:expert-advisory-corps} or other qualified body members exists from which to draw scoring panels, the \gls{gls:qualification-commons} bootstraps as follows: the \gls{gls:rubric-body} is initially populated by persons designated by the \gls{gls:constitutional-assembly}, serving one-time non-renewable founding terms, barred from serving in any body whose qualification pool they help populate (Article \ref{art:sortition-design}, \ref{art:sortition-design:s-qualification-commons}). Initial domain scoring panels are drawn reciprocally from adjacent domains rather than from within each domain, preventing any single professional community from controlling its own entry point at the founding moment. These founding arrangements sunset automatically once each domain pool reaches a threshold of one hundred fifty qualified members, sufficient to ensure that the probability of any individual member's selection in a given scoring cycle does not exceed five percent, at which point the standard \gls{gls:qualification-commons} process takes over.

\ArtSection{Cooperative Development Bank Bootstrap}{art:founding-mechanisms:s-cooperative-development-bank}
At founding, before a mature cooperative sector exists from which to draw Peer Panel members, lending decisions are made by panels drawn from the \gls{gls:constitutional-assembly}'s designees, with explicit sunset once the cooperative sector reaches one hundred established enterprises. The Development Partnership process and all other Bank architecture apply from day one.

\ArtSection{Land Tenure Conversion}{art:founding-mechanisms:s-land-tenure-conversion}
On territory where persons currently hold private land ownership at the time of ratification, existing titles are converted as follows:

Persons who constructed their dwelling or made improvements whose assessed labor value exceeds twenty-five percent of the assessed replacement cost receive a Residential Tenure automatically at ratification. No charge applies. The security, heritability, and non-transferability provisions of Article \ref{art:land-tenures}, \ref{art:land-tenures:s-residential-land} apply immediately.

Persons who purchased their dwelling under the prior system receive a Residential Tenure automatically at ratification on the same terms. In recognition of the financial investment made in good faith under the prior system, these persons are additionally eligible for a one-time Commons Transition Payment, administered by the \gls{gls:housing-implementation-council} against published criteria. The Commons Transition Payment is assessed at the difference between the replacement cost of improvements and the original purchase price adjusted for inflation, capped at the replacement cost of improvements. The payment phases out linearly for persons whose non-housing financial assets exceed the Commonwealth median annual wage, and is eliminated entirely for persons whose non-housing financial assets exceed ten times the Commonwealth median annual wage at the time of ratification. The Commonwealth median annual wage for this purpose is determined by the \gls{gls:monetary-authority} at ratification and published through the \gls{gls:transparency-engine}.

All existing tenants and renters under the prior system are immediately enrolled in Commons-Allocated Housing under Article \ref{art:open-presence}, \ref{art:open-presence:s-stable-housing-transition} and receive stable housing through that mechanism. No person is displaced by the transition.

Market transferability of all Residential Tenures ends at ratification. No Residential Tenure created under this Section may be sold, leased for profit, or used as a market commodity in any form after ratification.

Cooperatives that constructed or significantly developed their premises prior to ratification receive a Builder's Productive Tenure automatically at ratification. Cooperatives that purchased their premises under the prior system receive a Builder's Productive Tenure automatically at ratification and are additionally eligible for a one-time Commons Cooperative Transition Payment, assessed at the difference between improvement replacement cost and the inflation-adjusted original purchase price, capped at replacement cost. The payment is assessed against the cooperative's collective financial assets rather than individual member assets, and phases out linearly for cooperatives whose non-premises financial assets exceed ten times the Commonwealth median cooperative sector annual wage, eliminating entirely at the upper threshold. Market transferability of all Productive Tenures ends at ratification.

The \gls{gls:housing-implementation-council} administers all assessments under this Section against published criteria, with full documentation logged through the \gls{gls:transparency-engine}. Disputes are heard by the \gls{gls:housing-appeals-panel} under Article \ref{art:open-presence}, \ref{art:open-presence:s-commons-allocated-housing}. The \gls{gls:constitutional-court} (Article \ref{art:const-court}) has review authority on constitutional grounds.

\ArtSection{Staggered Sortition Bootstrap}{art:founding-mechanisms:s-staggered-sortition-bootstrap}
The staggered bootstrapping of founding sortition body populations -- including term-length lottery assignment, permanent non-renewability of founding service, and pro-rated civic service benefits -- is governed by Article \ref{art:sortition-design}, \ref{art:sortition-design:s-bootstrapping-first-legislature}.